In this section, we provide a new bound that relates the accuracy of Program~\ref{prog:primal} to the $p$-radius and the spectral radius of the adjacency matrix of the automaton.
The $p$-radius is a generalization of the joint spectral radius;
we recover the JSR with $p=\infty$.
The definition of the $p$-radius and its property can be found in \appref{pradius}; see \cite{jungers2011fast} for an introduction.

An important property of the $p$-radius is that it is increasing in $p$.
\begin{lemma}
  \label{lem:incp}
  Consider \defAG{}.
  For any integers $p \leq q$,
  \begin{equation}
    \label{eq:incp}
    \cpr[p] \leq \cpr[q] \leq \cjsr \leq \rho(A(\G))^{\frac{1}{q}}
    \cpr[q] \leq \rho(A(\G))^{\frac{1}{p}} \cpr[p].
  \end{equation}
\end{lemma}
The proof can be found in \appref{pradius}.
This Lemma is already known in the unconstrained case~\cite{zhou2002p}.

\begin{myrem}
  \lemref{incp} shows that the $p$-radius provides an upper and lower bound on the CJSR.
  It is known that the $2d$-radius can be computed
  either by computing a spectral radius~\cite{blondel2005computationally}
  or by solving a linear program~\cite{parrilo2008approximation}
  (see \cite{ogura2016efficient} for computation algorithms when $p$ is not an even integer).
  One needs to do either of these two operations on
  either the sum of the $d$th Kronecker power of the veronese 2-lift of the matrices
  or the sum of the veronese $2d$-lift of the matrices.
  However, the size of the matrices obtained by taking the $d$th Kronecker power grows rapidly
  and the veronese lifting needs to compute permanents which is computationally very demanding.
  Usually, for the same $d$, \progref{primal} gives approximations of the CJSR with a much higher accuracy at the cost of solving an SDP on smaller matrices.
  While solving an SDP is more demanding than computing a spectral radius or solving a linear program, since the matrices are smaller than with the Kronecker powering
  and we do not need to compute permanents, this method is actually more scalable.
\end{myrem}


In this section, we show the following bound stating that the solution found by \progref{primal} is at least as good as the bound obtained by computing the $2d$-radius.
\begin{theorem}
  \label{theo:second}
  Consider \defAG{} and a positive integer $d$.
  The approximation given by \progref{primal} using homogeneous polynomials of degree $2d$
  satisfies:
  \begin{equation}
    \label{eq:second}
    \jsrsos \leq \rho(A(\G))^{\frac{1}{2d}}\cdr \leq \rho(A(\G))^{\frac{1}{2d}} \cjsr
  \end{equation}
  where $A(\G)$ is the adjacency matrix of $\G$.
\end{theorem}
Note that the second inequality in \eqref{eq:second} is simply \eqref{eq:incp}.
This theorem is proved at the end of this section.

As a corollary of Theorem~\ref{theo:second}, in the trivial cases such that $\rho(A(\G)) = 1$, the approximation is exact.
This corresponds to the case where every node of $\G$ has indegree and outdegree 1.
In that case, the graph forms a cycle of some length $k$ and the CJSR is simply the $k$th root of the spectral radius of the product of the matrices along this cycle.

In the general case, the following approximation guarantee is known (note that the bound does not take into account the particular structure of the automaton):
\begin{theorem}[{\cite[Theorem~3.6]{philippe2016stability}}]
  \label{theo:matt}
  Consider \defAGn{} and a positive integer $d$.
  The approximation $\jsrsos$ given by Program~\ref{prog:primal} using homogeneous polynomials of degree $2d$
  satisfies:
  \[ \jsrsos \leq {n+d-1 \choose d}^{\frac{1}{2d}} \cjsr. \]
\end{theorem}

The results of \theoref{second}, \theoref{matt} and \eqref{eq:upperbound} are summarized by the following corollary.
\begin{corollary}
  \label{coro:mix}
  Consider \defAGn{} and a positive integer $d$,
  the approximation given by Program~\ref{prog:primal} using homogeneous polynomials of degree $2d$
  satisfies:
  \[
    \min\Big\{{n+d-1 \choose d}, \rho(A(\G))\Big\}^{-\frac{1}{2d}} \jsrsos%
    \leq \cjsr \leq \jsrsos
  \]
  where $A(G)$ is the adjacency matrix of the automaton $\G$.
\end{corollary}
We see that we can have arbitrary accuracy by increasing $d$.

For the arbitrary switching case, $\rho(A(\G))$ is equal to the number of matrices $m$.
Theorem~\ref{theo:second} was already known in this particular case \cite[Theorem~4.3]{parrilo2008approximation}.

Our proof technique relies on the analysis of an iteration in the vector space of polynomials of degree $2d$.
When this iteration converges, it converges to a feasible solution of Program~\ref{prog:primal}.
By analysing this iteration as affine iterations in this vector space, we derive a sufficient condition for its convergence
and thus an upper bound for $\jsrsos$.

Consider the iteration
\begin{align}
  \notag
  p_{v,0}(x) & = 0,\\
  \label{eq:iter}
  p_{v,k+1}(x) & = q_v(x) + \frac{1}{\iterc} %
  \sum_{(u,v,\sigma) \in \Arcs} p_{u,k}(A_\sigma x), \quad  v \in \Nodes
\end{align}
for fixed homogeneous polynomials $q_v(x)$ of degree $2d$ in $n$ variables (not necessarily different)
and a constant $\iterc > 0$.

When this iteration converges, it converges to a feasible solution of Program~\ref{prog:primal}.
\begin{lemma}
  \label{lem:conv}
  Consider a constant $\iterc > 0$.
  If there exist homogeneous polynomials $q_v(x)$ in the interior of the SOS cone such that
  iteration~\eqref{eq:iter} converges then
  \[
    \jsrsos \leq \iterc^{\frac{1}{2d}}.
  \]
  \begin{proof}
    Suppose the iteration converges to the polynomials $p_{v,\infty}(x)$.
    It is easy to show by induction that $p_{v,k}(x)$ is SOS for all $k$.
    It is trivial for $k = 0$ and if it is true for $k$ then it is also true for $k+1$ by \eqref{eq:iter}.
    Since the SOS cone is closed, $p_{v,\infty}$ is SOS.
    Now by \eqref{eq:iter}, for each $v \in \Nodes$,
    \[ p_{v,\infty}(x) = q_v(x) + \frac{1}{\iterc} %
    \sum_{(u,v,\sigma) \in \Arcs} p_{u,\infty}(A_\sigma x) \]
    so $p_{v,\infty}(x)$ is also in the interior of the SOS cone.
    For each edge $(u,v,\sigma)$, by manipulating the above equation, we have
    \[ \iterc p_{v,\infty}(x) - p_{u,\infty}(A_\sigma x) = \iterc q_v(x) + %
    \sum_{(u',v',\sigma') \in \Arcs} p_{u',\infty}(A_{\sigma'} x) \]
    so $\iterc p_{v,\infty}(x) - p_{u,\infty}(A_\sigma x)$ is SOS.
    Therefore $(\{\, p_{v,\infty}(x) : v \in \Nodes \,\}, \iterc^{\frac{1}{2d}})$ is a feasible solution of Program~\ref{prog:primal}.
  \end{proof}
\end{lemma}

In view of \lemref{conv}, it is thus natural to analyse
under which condition iteration~\ref{eq:iter} converges. Recall
that iteration~\ref{eq:iter} is an affine map on the vector space of
homogeneous polynomials of degree $2d$.

\begin{proof}[Proof of Theorem~\ref{theo:second}]
  Iteration~\ref{eq:iter} is an affine map on the vector space of
  homogeneous polynomials of degree $2d$.  It is well known that
  if the convergence is guaranteed when we only retain the linear
  part of the affine map then it is also guaranteed for the affine
  iteration.

  Therefore we can analyse instead the following iteration
  \begin{align*}
    p_{v,0}(x) & = q_v(x),\\
    p_{v,k+1}(x) & = \frac{1}{\iterc} %
    \sum_{(u,v,\sigma) \in \Arcs} p_{u,k}(A_\sigma x), \quad  v \in \Nodes
  \end{align*}

  We can see that
  \begin{equation}
    p_{v,k+1}(x) = \frac{1}{\iterc^k}\sum_{s \in \Arcsin_k(v)} q_{s(1)}(A_s x)\\
    \leq \frac{1}{\iterc^k}\sum_{s \in \Arcsin_k(v)} q_{s(1)}(A_s x).
  \end{equation}

  Consider a norm $\|\cdot\|$ of $\R^n$ and its corresponding induced matrix norm of $\R^{n \times n}$.
  For each $v \in \Nodes$, we know by continuity of $q_v(x)$ that there exist $\beta_v > 0$ such that
  \[ q_v(x) \leq \beta_v\|x\|^{2d} \]
  for all $x \in \R^n$.
  Let $\beta = \max_{v \in \Nodes} \beta_v$.

  Therefore,
  \begin{align*}
    p_{v,k+1}(x)
    & \leq \frac{1}{\iterc^k}\sum_{s \in \Arcsin_k(v)} \beta_v \|A_s\|^{2d} \|x\|^{2d}\\
    & \leq \frac{\beta}{\iterc^k} \|x\|^{2d} \sum_{s \in \Arcsin_k(v)} \|A_s\|^{2d}.
  \end{align*}
  By Lemma~\ref{lem:rhoAG} and Remark~\ref{rem:path}, if
  $\iterc > \rho(A(G))\cdr^{2d}$, $\lim_{k \to \infty} p_{v,k}(x) = 0$.
  We obtain the result by \lemref{conv}.
\end{proof}
